\documentclass{beamer}
\usepackage[utf8]{inputenc}

\usetheme{Berkeley}
\usecolortheme{default}

\setbeamertemplate{footline}[frame number]
\beamertemplatenavigationsymbolsempty

\setbeamerfont{subsection in toc}{size=\scriptsize}

\usepackage[most]{tcolorbox}

\title{Literature Managment}
\author[]{IB 514: Scientific Writing}
\date[]{}

\AtBeginSection[]
{
  \begin{frame}
    \frametitle{Overview}
    \tableofcontents[currentsection]
  \end{frame}
}



\begin{document}

%------------------------------------------------
\AtBeginSection[]
{
  \begin{frame}
    \frametitle{Overview}
    \tableofcontents[currentsection]
  \end{frame}
}

%------------------------------------------------
\begin{frame}
  \titlepage
\end{frame}

% ============================================================
\section{Framing the Challenge}
% ============================================================
\begin{frame}{The Literature Problem}

  \footnotesize
  \textbf{Two Competing Demands}
  \begin{itemize}
    \item Knowing the \textbf{foundational and classic} literature
      in your field
    \item Staying current with a \textbf{rapidly expanding} body
      of new work --- often across multiple subfields
  \end{itemize}

  \medskip
  \textbf{Common Failure Modes}
  \begin{itemize}
    \item \textbf{Passive accumulation:} PDF hoarding without
      reading or annotation
    \item \textbf{Over-reading without synthesis:} consuming
      papers without integrating them into an argument
    \item \textbf{Reliance on incomplete system:} searching
      only Google Scholar, relying only on specific keyword alert
  \end{itemize}

  \medskip
  \begin{block}{Key Principle}
    You do not need to read everything.
    You need a \textbf{structured system} for filtering,
    organising, and revisiting literature that balances skimming
    and deep reading.
  \end{block}

\end{frame}


% ============================================================
\section{Staying Current}
% ============================================================
% ------------------------------------------------------------
\subsection{Alerts and Feeds}
% ------------------------------------------------------------

\begin{frame}{Staying Current: Alerts and Feeds}

  \footnotesize
  \textbf{Set up alerts --- then let new literature come to you}

  \smallskip
  \begin{description}

    \item[\textbf{Journal ToC / bioR$\chi$v Alerts}]
      Subscribe through websites or aggregate
      via RSS reader (e.g., Feedly). Best for core journals in
      your subfield --- gives you everything, not just keyword matches.

    \item[\textbf{Web of Science / Google Scholar Alerts}]
      Key paper and keyword alerts delivered by email on a schedule
      you control. \textit{Highly filtered; low noise.}

    \item[\textbf{ResearchGate / Google Scholar Profile Suggestions}]
      Follow authors in your field, surface papers
      based on your reading history.

    \item[\textbf{Social Media}]
      Direct announcements of new papers; hashtags
      and community accounts surface work quickly.
      \textit{Informal but fast.}
  \end{description}

  \medskip
  \begin{block}{\textbf{Strategy}}
    Separate \textbf{scanning} from \textbf{deep reading}; schedule dedicated time for each.
  \end{block}

\end{frame}


% ------------------------------------------------------------
\subsection{Active Filtering}
% ------------------------------------------------------------

\begin{frame}{Staying Current: Active Filtering}

  \footnotesize
  \textbf{Three-Tier Filtering System}
  \smallskip
  \begin{enumerate}
    \item \textbf{Scan:} titles / abstracts only ---
      decide relevance in under 60 sec.
    \item \textbf{Triage:} save potentially relevant PDFs
      to reference manager or folder
    \item \textbf{Deep Read:} full reading reserved for
      highest-priority papers
  \end{enumerate}

  \medskip
  \textbf{Heuristics for Prioritisation}
  \begin{itemize}
    \item Direct relevance to current or future projects
    \item Methods, datasets, or frameworks of potential use
    \item Intriguing novelty / fun
  \end{itemize}

  \medskip
  \textbf{AI-Assisted Reading Tools (examples)}

  \smallskip
  \begin{tabular}{@{}lp{7.5cm}@{}}
    \hline
    NotebookLM  & Query PDFs conversationally;
                           rapid orientation to new subfield \\
    Speechify   & Convert papers to audio for
                           assessment in low-focus time \\
    \hline
  \end{tabular}

  \vfill
  \begin{block}{Reading everything is inefficient}
    \textbf{Selective re-reading} is where understanding develops.
  \end{block}

\end{frame}


% ============================================================
\section{Literature Search}
% ============================================================
% ------------------------------------------------------------
\subsection{Search Strategies}
% ------------------------------------------------------------
\begin{frame}{Literature Search Strategies (Traditional)}

  \footnotesize
  \textbf{Citation Chaining}
  \begin{itemize}
    \item \textbf{Backward search:} scan for common references in recent papers
    \item \textbf{Forward search:} scan and follow citations of key papers
  \end{itemize}

  \medskip
  \textbf{Author Search}
  \begin{itemize}
    \item Identify 5--10 leading (frequently cited) researchers;
    track the citation records of most cited papers
  \end{itemize}

  \medskip
  \textbf{Start Broad, Then Narrow}
  \begin{itemize}
    \item Use general keywords to identify core papers;
      refine iteratively using Boolean operators
      (\texttt{AND}, \texttt{OR}, \texttt{NOT})
    \item Update as terminology, useful keywords and authors become clearer
  \end{itemize}

  \medskip
  \begin{block}{Goal (especially for new-to-you fields)}
    Identify the \textbf{core intellectual structure} of a topic,
    not just a list of papers.
  \end{block}

\end{frame}


% ------------------------------------------------------------
\subsection{Search Tools}
% ------------------------------------------------------------
\begin{frame}{Literature Search Tools}

  \footnotesize
  \begin{description}
    \item[\textbf{Web of Science}]
      Curated index; strong citation tracking and bibliometric
      tools; advanced Boolean searches and exportable
      results. The Standard for systematic searches.
      {\scriptsize \textit{Requires institutional access.}}

    \item[\textbf{Google Scholar}]
      Far less precise than WoS but broad coverage including preprints, theses, and
      grey literature. 
      Useful for quick searches and obtaining PDFs.
      {\scriptsize \textit{Free.}}

    \item[\textbf{AI-Assisted Tools}]
       {\scriptsize Tools evolving rapidly. Reliability growing but TBD.}\\
      \begin{itemize}\footnotesize
        \item \textbf{Elicit:} structured queries
          and paper summaries
        \item \textbf{Scite:} identifies whether papers \emph{support} or
          \emph{contrast} the work they cite
        \item \textbf{Semantic Scholar / ResearchRabbit:}
          relevance ranking, topic mapping, and citation
          network visualisation
      \end{itemize}
  \end{description}

\end{frame}


% ============================================================
\section{Reference Management}
% ============================================================
% ------------------------------------------------------------
\subsection{Tools}
% ------------------------------------------------------------
\begin{frame}{Reference Managers}

  \footnotesize

  \begin{block}{\textbf{Core Function}}
    Database to organize sources and auto-format citations and bibliographies
  \end{block}
  \textbf{Other Functions}
  \begin{itemize}\scriptsize
    \item Enable tagging, search, and annotation
    \item Store and organise PDFs
    \item Sync across devices; shared group libraries
  \end{itemize}

  \medskip
  \textbf{Zotero\footnote{{\tiny Free (maybe limited) and open-source; $\phantom{}^2$\$\$\$ proprietary}}} (most used?)
  \begin{itemize}
     \item Integrates with Word, Google, \LaTeX{}
      (w/ \texttt{Better BibTeX})
    \item Browser plugin captures
      metadata and PDFs in one click
    \item Group libraries: share collections with
      co-authors in real time
    \item Supports tags, notes, related items, PDF annotation
  \end{itemize}

  \medskip
  \scriptsize
  \textbf{Alternatives include:} 
  BibDesk$\phantom{}^1$ (\LaTeX), 
  Mendeley$\phantom{}^2$ (strong PDF reader),
  EndNote$\phantom{}^2$ (institutional standard), 
  Paperpile$\phantom{}^2$ (Google Docs integration)

\end{frame}


% ------------------------------------------------------------
\subsection{Organisation and Annotation}
% ------------------------------------------------------------
\begin{frame}{Organizing Sources}

  \footnotesize
  \textbf{Build Structure, Not Just Storage}
  \begin{itemize}
    \item Organise by \textbf{project} or \textbf{manuscript},
      not by topic alone --- projects have deliverables
    \item Use \textbf{tags} for cross-cutting themes
      {\scriptsize (e.g., \texttt{classic},
      \texttt{2read}, \texttt{cited-in-ch2})}
  \end{itemize}

  \medskip
  \textbf{Recommended Approach: Two-Layer System}
  \begin{itemize}
    \item \textbf{Reference manager:} storage, PDF access, citation generation
    \item \textbf{External notes} (e.g., a plain-text or
      Markdown file per project): conceptual synthesis,
      argument mapping, connections across papers
  \end{itemize}
  \begin{block}{The two layers complement each other}
    The reference manager stores, the notes file thinks
  \end{block}
\end{frame}


% ------------------------------------------------------------
\begin{frame}{Annotated Bibliographies}

  \footnotesize
  \textbf{Purpose}
  \begin{itemize}
    \item Convert reading into structured, reusable knowledge
    \item Force active engagement
    {\scriptsize (you can't annotate what you've only skimmed)}
    \item Becomes raw material for your Introduction and
      Discussion sections
  \end{itemize}

  \medskip
  \textbf{What to Include per Entry}
  \begin{itemize}
    \item Core question or objective of the paper
    \item Key findings and conclusions
    \item Methods or analytical approach
    \item Relevance to your own work
    \item Limitations or unresolved gaps
  \end{itemize}

  \medskip
  \textbf{Suggestions}
  \begin{itemize}
    \item Write in your own words --- not copy-pasted from abstract
    \item Keep entries concise: 3--6 bullet points
    \item Use Zotero's note field to keep with reference
    \item Update entries as your understanding evolves
  \end{itemize}

\end{frame}


% ============================================================
\section{Writing from the Literature}
% ============================================================
\subsection{Integration Strategies}
% ------------------------------------------------------------
\begin{frame}{Integrating Literature in Your Writing}

  \footnotesize

  \textbf{Writing \textit{with} specific papers in mind}
  \begin{itemize}
    \item \emph{Pros:} precise citation; accurate representation
      of prior work; anchors claims to evidence
    \item \emph{Cons:} can produce fragmented,
      citation-heavy prose; limits synthesis; writing
      follows the literature rather than your argument
    \item \emph{Best for:} methods sections, direct
      comparisons, specific statistics
  \end{itemize}

  \medskip
  \textbf{Writing \textit{without} specific papers in mind}
  \begin{itemize}
    \item \emph{Pros:} promotes conceptual clarity;
      argument-driven structure; encourages synthesis, faster drafting%
      {\scriptsize \footnote{\textit{"Write drunk, edit sober."} --- Hemingway}}
    \item \emph{Cons:} risk of vague/unsupportable claims;
      requires well-stocked library; misrepresented citations
    \item \emph{Best for:} introduction and discussion (synthesis)
  \end{itemize}

\end{frame}


% ------------------------------------------------------------
\subsection{Hybrid Writing Strategy}
% ------------------------------------------------------------

\begin{frame}{Best Practice: The Hybrid Writing Strategy}

  \footnotesize
  \textbf{Step 1 --- Draft Conceptually}
  \begin{itemize}
    \item Write the argument without citations
    \item Focus on logic, structure, and your own voice
    \item Placeholder (\texttt{[CITE]}) where citation is needed
  \end{itemize}

  \medskip
  \textbf{Step 2 --- Map Literature Onto the Argument}
  \begin{itemize}
    \item Add citations to support each claim
    \item Replace general statements with specific evidence
    \item Consult annotated bibliography ---
      most sources are likely there
  \end{itemize}

  \medskip
  \textbf{Step 3 --- Refine for Balance}
  \begin{itemize}
    \item Avoid over-citation --- not every sentence
      needs a reference
    \item Ensure key and foundational works are represented
    \item Verify every citation against the source
    \item Never cite a paper you have not fully read
  \end{itemize}

  \medskip
  \begin{block}{Recommendation}
    Use argument-first writing for structure and flow,
    then verify every citation.
  \end{block}

\end{frame}



% ============================================================
\section{Takeaway}
% ============================================================
\begin{frame}{Takeaway}

  \footnotesize

  Knowing the classics and staying on top of the literature is a \textbf{process} 
  and is a part of your \textbf{job} as a scientist (and educator).

  \smallskip
  Recognize it as such and \textbf{regularize it} as part of your weekly schedule.

  \medskip
  \textbf{Remember the four pillars:}
  \begin{description}
     \item[\textbf{Discover:}] alerts + citation chaining + active search 
    \item[\textbf{Filter:}] scan $\rightarrow$ triage $\rightarrow$ deep read
    \item[\textbf{Organise:}] reference manager + external synthesis notes + annotated bibliography
    \item[\textbf{Integrate:}] argument-first drafting $\rightarrow$ citation mapping $\rightarrow$ refinement
  \end{description}

  \medskip
  \begin{block}{Takeaway}
    Managing the literature is not about volume.\\
    It's about \textbf{structure, selectivity, and synthesis}.
  \end{block}

\end{frame}

\end{document}